\chapter{Experiments and exercises}
\label{ch:exercises}

\section{Experiments with the program}
These experiments need no programming. Each one connects something visible in the program to a
section of Part~I.

\begin{enumerate}[label=\textbf{E\arabic*.}]
\item \emph{Texture space} (Chapter~\ref{ch:texspace}). Select the ``UV coordinates'' view. On each
solid, locate $u=0$, $u=1$ and the seam; on the sphere, observe the pinching at the poles; on the
cylinder, explain why the caps show the same colours as parts of the side.
\item \emph{Normalized coordinates.} Hover over the corners of the large image in the Inspector and
check that $(u,v)$ goes from $(0,0)$ at the top left to $(1,1)$ at the bottom right, and that the
texel index is $\lfloor uW\rfloor$.
\item \emph{Bump versus displacement} (Chapter~\ref{ch:height}). With the sphere and a rough
material, untick ``Bump'' and look at the silhouette and the lighting. Then tick ``Bump'' and
untick ``Displacement''. Then increase ``Displacement scale'' to $0.3$ with bump off: why does the
surface look strangely smooth despite its jagged outline?
\item \emph{Cracks.} With displacement on, look at the edges of the pyramid and the rims of the
cylinder from close up. Explain the cracks with equation~\eqref{eq:displace}.
\item \emph{Bump strength} (Section~\ref{sec:bumpformula}). Pin the probe on a groove and read
$\nts$ while varying the bump strength from $0$ to $30$. Plot the angle between $\nts$ and
$(0,0,1)$ as a function of $k$ and compare with $\arctan\bigl(k\sqrt{h_x^2+h_y^2}\bigr)$.
\item \emph{Tangent space versus world space} (Chapter~\ref{ch:tangent}). Pin the probe and start
the rotation. Which numbers change and which do not? Then change the shape: which numbers change
now?
\item \emph{Normal maps} (Section~\ref{sec:normalmaps}). For several texels, compare the signs of
the $x$ and $y$ components of the computed $\nts$ and of the pack normal map. Do they agree? What
would happen with an OpenGL normal map?
\item \emph{Roughness} (Chapter~\ref{ch:lighting}). Untick ``Roughness'' and sweep the constant
roughness from $0.04$ to $1$ while the light orbits automatically. Relate the size of the
highlight to Figure~\ref{fig:ggx}.
\item \emph{Colour spaces} (Section~\ref{sec:srgb}). Compare the Inspector's ``Color (sRGB)'' value
of a mid-grey texel with the value that the lighting actually uses after decoding. How large is
the difference?
\end{enumerate}

\section{Programming exercises}
In increasing order of difficulty. After editing a shader, press F5; after editing C++,
rebuild.

\begin{enumerate}[label=\textbf{P\arabic*.}]
\item In \code{SpherePatches}, remove the minus signs of the $z$ components. Describe what
happens to the texture and to the Inspector's $\Tv$ and $\Bv$, and relate it to
equation~\eqref{eq:orientation} and to the sign $w$.
\item Replace the central differences by forward differences, both in \code{BumpNormalTS} and in
\code{BumpNormalTangent}. Compare the probe's normal before and after on a sharp edge.
\item Add a view mode that shows the angle between the bumped normal and the geometric normal as a
heat map.
\item Add a view mode that shows the difference between the normal computed from the bump map and
the material's normal map. (You will need to bind the normal map to a fifth texture slot.)
\item Make the bump strength physically correct by dividing by the local stretch of the
parametrization, $\lVert\vect P_u\rVert$ and $\lVert\vect P_v\rVert$
(Section~\ref{sec:bumpformula}). On which solid is the change most visible?
\item Recompute the normal in the vertex shader from the displacement map, so that displacement
alone produces correct shading.
\item Close the cracks of the pyramid by displacing the vertices of the edges along the average of
the normals of the two faces that meet there.
\item Replace vertex displacement by hardware tessellation (hull and domain shaders) with a
tessellation factor that depends on the distance to the camera.
\end{enumerate}
